\providecommand{\toplevelprefix}{../..}  %
\documentclass[../../book-main.tex]{subfiles}

\begin{document}

\chapter*{Prefacio}
% \vspace{-2cm}
\begin{center}
    «{\em Todos los caminos conducen a Roma}.»

\end{center}
\vspace{5mm}

\paragraph{Objetivo.}
Este libro revela y estudia un problema común y fundamental que subyace
a casi todas las prácticas modernas de inteligencia artificial.
Se trata de {\em ¿cómo aprender de manera eficaz y eficiente una
    distribución de datos de baja dimensión en un espacio de alta dimensión y
    luego transformar esta distribución en una representación compacta y 
estructurada?} Para cualquier sistema inteligente, natural o artificial,
 una representación de este tipo puede considerarse en general como {\em memoria} o
{\em conocimiento} (empírico) aprendido a partir de datos percibidos del mundo externo. 
En los últimos años, a menudo se le denomina informalmente un «modelo del mundo.»

\paragraph{Público objetivo.}
Este libro de texto tiene como objetivo proporcionar una introducción sistemática a los
principios matemáticos y computacionales del aprendizaje de representaciones 
(profundas) de tales distribuciones de datos, como forma computable de
memoria, para {\em estudiantes de último año de licenciatura y de posgrado principiantes}.
Los principales prerequisitos para este libro son el álgebra lineal a nivel universitario,
probabilidad y estadística, así como optimización. Cierta familiaridad con 
los conceptos básicos del procesamiento de señales (representaciones dispersas y detección
comprimida en particular), de teoría de la información, así como de control por
retroalimentación mejoraría la comprensión del lector.

\paragraph{Motivación.}
La principal motivación para escribir este libro es que ha habido
tremendos avances en estos últimos años, tanto por los autores como por
numerosos colegas, que buscan establecer un enfoque riguroso y basado en
principios para entender las redes neuronales profundas y, más en general, la
inteligencia como tal. La metodología deductiva defendida por este nuevo 
enfoque, que es altamente complementaria, está en contraste directo con la 
metodología dominante que subyace a las prácticas contemporáneas en inteligencia 
artificial, que son inductivas y basadas en ensayo y error. La falta de comprensión
de esos poderosos modelos y sistemas de IA ha generado un creciente
entusiasmo y temor en el seno de la sociedad. Creemos que es necesario, más 
que nunca, hacer un esfuerzo serio por establecer un enfoque basado en principios 
para comprender lo que es la inteligencia. Un objetivo general de este libro es 
proporcionar pruebas teóricas y experimentales sólidas que demuestren que es ahora
posible estudiar la inteligencia como disciplina científica y matemática. Argumentaremos
que la inteligencia es la capacidad fundamental de desarrollar nuevas memorias (o
conocimientos) o de corregir las existentes. Por lo tanto, podemos considerar este
libro como un primer intento de desarrollar {\em una teoría matemática de la 
inteligencia}, al nivel de aprender conocimientos empíricos en forma de memoria, 
tal como sugiere el subtítulo del libro.

A nivel técnico, el marco teórico presentado en este libro ayuda a 
reconciliar una brecha antigua entre el enfoque clásico del
modelado de estructuras de datos principalmente basado en
modelos analíticos geométricos, algebráicos y probabilísticos 
(p.ej. subespacios, distribuciones gaussianas y ecuaciones) y el enfoque «moderno»
basado en modelos no paramétricos diseñados empíricamente a partir de datos
(p.ej. redes profundas). Como resultado, la unificación de
estas dos metodologías aparentemente separadas es posible e
incluso natural si se entiende que todas tratan de aprender y
representar estructuras {\em de baja dimensión} en la
distribución de datos de interés. Son simplemente maneras
diferentes de descubrir, representar y explotar estas estructuras 
de baja dimensión. Desde esta perspectiva, incluso muchas
técnicas computacionales aparentemente no relacionadas,
desarrolladas independientemente en campos separados y en distintos
momentos, pueden ser comprendidas mejor ahora en un marco teórico
y computacional común y es probable que puedan estudiarse en su
conjunto a partir de ahora. Como veremos en este libro, estas 
técnicas, entre otras, incluyen: 
\begin{enumerate}
    \item La codificación y decodificación compresiva con pérdidas
        desarrollada en la teoría de la información clásica y la
        teoría de la codificación frente al aprendizaje de 
        representaciones moderno;
    \item La reducción de ruido y la reducción de dimensionalidad
        en el procesamiento clásico de señales frente a los modelos
        de difusión y reducción de ruido de los métodos generativos
        modernos;
    \item La inferencia bayesiana por vía del máximo a posteriori o
        el muestreo condicionado frente a la optimización con 
        restricciones por la vía de técnicas de continuación
        \footnote{como el método del lagrangiano aumentado}.
\end{enumerate}

\paragraph{Contenido principal.}
Creemos que el marco conceptual y computacional unificado que se
presenta en este libro será de gran utilidad para aquellos lectores
que verdaderamente quieran aclarar los misterios y malentendidos
de las redes neuronales profundas y la inteligencia (artificial).
Más allá de lo anterior, este marco pretende equipar a los lectores 
con principios rectores para desarrollar sistemas significativamente
mejores y verdaderamente inteligentes en el futuro. Más específicamente,
aparte de una introducción (capítulo) informal, el contenido técnico
principal del libro se organizará en forma de seis temas (capítulos)
íntimamente relacionados:
\begin{enumerate}
    \item Comenzaremos con los modelos clásicos y más básicos del
        análisis de componentes principales (PCA), del análisis de
        componentes independientes (ICA) y del aprendizaje de 
        diccionarios (DL), las que suponen que las distribuciones
        de baja dimensión de interés poseen estructuras
        lineales e independientes. A partir de estos modelos simples
        idealizados que han sido ampliamente estudiados y comprendidos
        en procesamiento de señales y detección por compresión, 
        introduciremos las ideas más fundamentales\footnote{tales 
        como a reducción de ruidos y la reducción de dimensionalidad} sobre
        cómo aprender y cómo representar distribuciones de baja
        dimensión.

    \item Para generalizar estos modelos clásicos y extender
        sus soluciones a distribuciones de baja dimensión que no
        puedan ser representadas fácilmente por un modelo analítico
        sencillo, introducimos un principio computacional universal para
        el aprendizaje de tales distribuciones: {\em la compresión}.
        Como veremos, la compresión de datos proporciona una visión 
        unificadora de todos los enfoques clásicos y modernos aparentemente
        diferentes para el aprendizaje de distribuciones o representaciones,
        incluyendo la minimización de entropía, la reducción de ruido por
        basado en score-matching, compresión con pérdidas y distorsión de 
        tasa, así como el aprendizaje discriminativo de representaciones
        por vía de maximizar la ganancia de información.

    \item En este marco unificador, las redes neuronales profundas (DNNs),
        tales como ResNet, CNN y transformadores, pueden ser todas 
        interpretadas matemáticamente como algoritmos de optimización
        (desenrollados) que iterativamente logran mejor compresión
        y mejor representación al reducir el largo o la tasa de la 
        codificación o al ganar información. Este marco no solo ayuda
        a explicar las arquitecturas de redes profundas empíricamente
        diseñadas hasta ahora, sino que también conduce a nuevos
        diseños arquitectónicos que pueden ser significativamente 
        más simples y más eficientes.

    \item Más allá de lo anterior, a fin de asegurar que la representación
        aprendida para una distribución de datos sea correcta y consistente,
        se hacen necesarias las arquitecturas de {\em autocodificadores}
        que incluyen tanto codificadores como decodificadores. Para que
        un sistema de aprendizaje sea plenamente automático y de operación
        contínua, introduciremos un poderoso {\em esquema de transcripción
        en circuito cerrado} que permite quee un sistema autocodificador se
        autocorrija y, por tanto, se autoperfeccione, por medio
        de un juego minimax entre codificador y decodificador.

    \item Para conectar la teoría con la práctica, luego estudiaremos cómo
        la distribución de datos y su representación aprendidas pueden
        utilizarse como un potente a priori para llevar a cabo inferencias
        bayesianas u optimizaciones con restricciones que se manifiestan
        en casi todos los tipos de tareas y configuraciones más ampliamente utilizadas en
        la práctica de la inteligencia artificial moderna, incluyendo la 
        estimación, el completamiento y la generación condicionales de datos
        reales de alta dimensión, tales como imágenes y textos.     

    \item Por último, pero no menos importante, mostraremos paso a paso cómo 
        aprender de manera efectiva y eficiente las representaciones profundas
        de distribuciones de baja dimensión a partir de conjuntos
        de datos del mundo real a gran escala, incluyendo datos visuales, movimientos 
        corporales y datos textuales, usándolos en numerosas aplicaciones
        prácticas, tales como clasificación, completado, segmentación, 
        generación de imágenes, estimación de movimiento y tareas similares
        para datos de texto.
\end{enumerate}

\paragraph{Resumen.}
Para resumir, el contenido técnico presentado en este libro
establece fuertes conexiones conceptuales y técnicas entre
el enfoque analítico clásico y el enfoque moderno impulsado por
datos, entre modelos paramétricos simples y los modelos
no paramétricos profundos, entre las diversas prácticas inductivas
y un enfoque deductivo unificado basado en principios básicos.
Revelaremos que numerosos enfoques aparentemente no relacionados
o incluso contradictorios entre sí, aunque se hayan desarrollado en 
campos distintos, con terminologías diversas y en momentos históricos diferentes,
sin embargo todos aspiran a lograr un objetivo común:
\begin{quote}
    \centering{\em perseguir y explotar las estructuras intrínsecas
    de baja dimensión de distribuciones de datos incrustadas
    en espacios de alta dimensión.}
\end{quote}
A tal fin, el libro nos guiará a través de un recorrido completo desde
la formulación teórica a la deducción matemática y luego a la 
implementación computacional y las aplicaciones prácticas.

\end{document}
